% Options for packages loaded elsewhere
\PassOptionsToPackage{unicode}{hyperref}
\PassOptionsToPackage{hyphens}{url}
\documentclass[
  10pt,
]{article}
\usepackage{xcolor}
\usepackage[margin=1in]{geometry}
\usepackage{amsmath,amssymb}
\setcounter{secnumdepth}{5}
\usepackage{iftex}
\ifPDFTeX
  \usepackage[T1]{fontenc}
  \usepackage[utf8]{inputenc}
  \usepackage{textcomp} % provide euro and other symbols
\else % if luatex or xetex
  \usepackage{unicode-math} % this also loads fontspec
  \defaultfontfeatures{Scale=MatchLowercase}
  \defaultfontfeatures[\rmfamily]{Ligatures=TeX,Scale=1}
\fi
\usepackage{lmodern}
\ifPDFTeX\else
  % xetex/luatex font selection
\fi
% Use upquote if available, for straight quotes in verbatim environments
\IfFileExists{upquote.sty}{\usepackage{upquote}}{}
\IfFileExists{microtype.sty}{% use microtype if available
  \usepackage[]{microtype}
  \UseMicrotypeSet[protrusion]{basicmath} % disable protrusion for tt fonts
}{}
\makeatletter
\@ifundefined{KOMAClassName}{% if non-KOMA class
  \IfFileExists{parskip.sty}{%
    \usepackage{parskip}
  }{% else
    \setlength{\parindent}{0pt}
    \setlength{\parskip}{6pt plus 2pt minus 1pt}}
}{% if KOMA class
  \KOMAoptions{parskip=half}}
\makeatother
\usepackage{graphicx}
\makeatletter
\newsavebox\pandoc@box
\newcommand*\pandocbounded[1]{% scales image to fit in text height/width
  \sbox\pandoc@box{#1}%
  \Gscale@div\@tempa{\textheight}{\dimexpr\ht\pandoc@box+\dp\pandoc@box\relax}%
  \Gscale@div\@tempb{\linewidth}{\wd\pandoc@box}%
  \ifdim\@tempb\p@<\@tempa\p@\let\@tempa\@tempb\fi% select the smaller of both
  \ifdim\@tempa\p@<\p@\scalebox{\@tempa}{\usebox\pandoc@box}%
  \else\usebox{\pandoc@box}%
  \fi%
}
% Set default figure placement to htbp
\def\fps@figure{htbp}
\makeatother
\setlength{\emergencystretch}{3em} % prevent overfull lines
\providecommand{\tightlist}{%
  \setlength{\itemsep}{0pt}\setlength{\parskip}{0pt}}
% _latex_preamble.tex
\usepackage{booktabs}
\usepackage{longtable}
\usepackage{geometry}
\usepackage{array}
\renewcommand{\arraystretch}{1.15}
\geometry{margin=1in}
\usepackage{float}
\usepackage{caption}
\usepackage{fancyhdr}
\pagestyle{fancy}
\fancyhf{}
\rhead{\thepage}
\lhead{Regional Skew Covariates}
\usepackage{tocloft}
\renewcommand{\cftsecfont}{\normalfont}
\renewcommand{\cftsecpagefont}{\normalfont}
\usepackage{booktabs}
\usepackage{longtable}
\usepackage{array}
\usepackage{multirow}
\usepackage{wrapfig}
\usepackage{float}
\usepackage{colortbl}
\usepackage{pdflscape}
\usepackage{tabu}
\usepackage{threeparttable}
\usepackage{threeparttablex}
\usepackage[normalem]{ulem}
\usepackage{makecell}
\usepackage{xcolor}
\usepackage{bookmark}
\IfFileExists{xurl.sty}{\usepackage{xurl}}{} % add URL line breaks if available
\urlstyle{same}
\hypersetup{
  pdftitle={Data Dictionary: Covariates by Domain},
  pdfauthor={C.J. Tinant with ChatGPT 4o (20250504)},
  hidelinks,
  pdfcreator={LaTeX via pandoc}}

\title{Data Dictionary: Covariates by Domain}
\author{C.J. Tinant with ChatGPT 4o (20250504)}
\date{}

\begin{document}
\maketitle

{
\setcounter{tocdepth}{2}
\tableofcontents
}
\section*{Introduction}
\addcontentsline{toc}{section}{Introduction}

This document summarizes the covariate metadata used in the regional
skew estimation study. Variables are grouped by \textbf{Domain} and
described in terms of their name, definition, unit, spatial scale, and
analytical notes.

\section*{ Climate }
\addcontentsline{toc}{section}{ Climate }

\begingroup\fontsize{8}{10}\selectfont

\begin{longtable}[t]{lp{8cm}l}
\caption{\label{tab:create_tables}Table A: Climate — Definition and Units}\\
\toprule
Short Name & Description & Unit\\
\midrule
\endfirsthead
\caption[]{Table A: Climate — Definition and Units \textit{(continued)}}\\
\toprule
Short Name & Description & Unit\\
\midrule
\endhead

\endfoot
\bottomrule
\endlastfoot
Longitude & Point-based (gage-level) data used as a station-level predictors, not
derived from zonal aggregation. & decimal degrees\\
Latitude & Point-based (gage-level) data used as a station-level predictors, not
derived from zonal aggregation. & decimal degrees\\
Climate Zone & Dominant Koppen Geiger climate class, aggregated to macroregional
zones. & categorical code\\
PHZM Zone Count & Count of unique USDA Plant Hardiness Zones, aggregated to
macroregional zones. & count\\
PHZM Zone (Dominant) & Dominant USDA Plant Hardiness Zone, aggregated to macroregional zones. & categorical code\\
\addlinespace
Annual Precip (Median) & Zonal median of total annual precipitation from PRISM 1991 to 2020
climatological normals, aggregated to Level II ecoregions. & mm\\
Annual Temp (Median) & Zonal median of total annual precipitation from PRISM 1991 to 2020
climatological normals, aggregated to Level II ecoregions. & degrees C\\
Pct Annual Precip May to Aug (Mean) & Zonal mean of the percentage of annual precipitation falling May to
August from PRISM 1991 to 2020 climatological normals, aggregated to
Level II ecoregions. & proportion [0 to 1]\\
Pct Annual Precip May to Aug (Median) & Zonal mean of the percentage of annual precipitation falling May to
August from PRISM 1991 to 2020 climatological normals, aggregated to
Level II ecoregions. & proportion [0 to 1]\\
Seasonal Precip StDev & Standard deviation of seasonal precip totals from PRISM 1991 to 2020
normals aggegated to Level III ecoregions. & mm\\
\addlinespace
Seasonal Precip IQR & IQR (25th to 75th percentile) of seasonal precip totals from PRISM
1991 to 2020 normals aggegated to Level III ecoregions. & mm\\
Fall Precip (Median) & Sum of zonal medians of fall monthly precip totals (Sep to Nov) from
PRISM 1991 to 2020 normals aggegated to Level III ecoregions. & mm\\
Winter Precip (Median) & Sum of zonal medians of winter monthly precip (Dec to Feb) from PRISM
1991 to 2020 normals aggegated to Level III ecoregions. & mm\\
Spring Precip (Median) & Sum of zonal medians of summer monthly precip (Jun to Aug) from PRISM
1991 to 2020 normals aggegated to Level III ecoregions. & mm\\
Summer Precip (Median) & Sum of zonal medians of spring monthly precip (Mar to May) from PRISM
1990 to 2020 normals aggegated to Level III ecoregions. & mm\\
\addlinespace
Freeze-Thaw Days & Zonal count of freeze-thaw transition days from PRISM daily data (1991
to 2020), aggegated to NHD+ catchments. & days\\
Precip Intensity (95th-percentile) & Zonal median of 95th percentile of daily precip on wet days from PRISM
daily normals (1991 to 2020), aggegated to NHD+ catchments. & mm\\
Wet Day Frequency (Seasonal) & Zonal median number of wet days (less than 0.1 mm/day) per season from
PRISM daily normals (1991 to 2020), aggregated to NHD+ catchments. & days\\*
\end{longtable}
\endgroup{}

\vspace{0.5cm}

\begingroup\fontsize{8}{10}\selectfont

\begin{longtable}[t]{lp{8cm}l}
\caption{\label{tab:create_tables}Table B: Climate — Notes and Temporal Scale}\\
\toprule
Short Name & Notes & Scale\\
\midrule
\endfirsthead
\caption[]{Table B: Climate — Notes and Temporal Scale \textit{(continued)}}\\
\toprule
Short Name & Notes & Scale\\
\midrule
\endhead

\endfoot
\bottomrule
\endlastfoot
Longitude & Represents site location in degrees east or west. Useful for detecting
climatic or physiographic gradients (moisture availability, storm
type). & Station\\
Latitude & Represents site location in degrees north. Commonly used to capture
latitudinal patterns in climate and hydrology (e.g., temperature,
seasonality). & Station\\
Climate Zone & Provides broad-scale climate context (e.g., humid continental vs.
semi-arid). Useful for regional stratification. & Macroregional\\
PHZM Zone Count & Reflects cold hardiness diversity; transitional zones may influence
vegetation, snowmelt timing, and freeze-thaw cycles. & Macroregional\\
PHZM Zone (Dominant) & Captures dominant cold tolerance regime based on minimum winter
temperatures. Relevant for snowmelt timing and freeze-thaw
sensitivity. & Macroregional\\
\addlinespace
Annual Precip (Median) & Represents typical annual moisture supply; important for soil moisture
availability, runoff potential, and vegetation productivity. & Regional\\
Annual Temp (Median) & Governs evapotranspiration, freeze thaw, and snowpack formation.
Represents the baseline thermal regime. & Regional\\
Pct Annual Precip May to Aug (Mean) & Characterizes warm-season water input. Useful for identifying
storm-driven systems tied to convective seasonality. & Regional\\
Pct Annual Precip May to Aug (Median) & Captures the typical proportion of precipitation falling in the warm
season, which is important for Great Plains hydrology. & Regional\\
Seasonal Precip StDev & Captures variability across seasons. Sensitive to outliers—useful for
detecting flashiness or hydrologic irregularity. & Subregional\\
\addlinespace
Seasonal Precip IQR & Characterizes intra-annual precipitation variability; robust to
outliers. & Subregional\\
Fall Precip (Median) & Captures fall precipitation; relevant to baseflow reset and transition
to frontal storm systems. & Subregional\\
Winter Precip (Median) & Indicates snowpack formation potential and winter water storage. & Subregional\\
Spring Precip (Median) & Key springtime recharge and flood generation season; relevant to
melt-transition systems and early crop needs. & Subregional\\
Summer Precip (Median) & Captures convective storm seasonality (flood-prone regions). & Subregional\\
\addlinespace
Freeze-Thaw Days & Captures thermally transitional zones; relevant for soil structure,
infiltration disruption, and freeze thaw runoff events. & Local\\
Precip Intensity (95th-percentile) & Captures storm severity at the upper tail. Correlated with flash flood
potential and runoff generation efficiency. & Local\\
Wet Day Frequency (Seasonal) & Describes storm frequency and seasonal spacing. Interacts with
intensity to shape hydrograph form and flood likelihood. & Local\\*
\end{longtable}
\endgroup{}

\clearpage

\section*{ Topography }
\addcontentsline{toc}{section}{ Topography }

\begingroup\fontsize{8}{10}\selectfont

\begin{longtable}[t]{lp{8cm}l}
\caption{\label{tab:create_tables}Table A: Topography — Definition and Units}\\
\toprule
Short Name & Description & Unit\\
\midrule
\endfirsthead
\caption[]{Table A: Topography — Definition and Units \textit{(continued)}}\\
\toprule
Short Name & Description & Unit\\
\midrule
\endhead

\endfoot
\bottomrule
\endlastfoot
Station Altitude & Point-based (gage-level) data used as a station-level predictors, not
derived from zonal aggregation. & meters\\
Mean Slope & Mean slope (from NED), aggregated to macroregional zones. & degrees\\
Median Slope & Median slope (from NED), aggregated to macroregional zones. & degrees\\
Altitude Zone & Mean elevation of gage locations, aggregated to macroregional zones. & meters\\
Mean Slope Distribution & Zonal mean of NED slope values, aggregated to Level II ecoregions. & degrees\\
\addlinespace
Slope Skewness & Zonal skewness of NED slope values, aggregated to Level II ecoregions. & unitless (skewness)\\
Slope Variability (StDev) & Zonal standard deviation (SD) of slope values from NED data,
aggregated to Level II ecoregions. & degrees\\
Slope Variability (IQR) & Zonal interquartile range (IQR) of NED slope values, aggregated to
Level II ecoregions. & degrees\\
Surface Roughness & Zonal median of surface roughness, derived from 10 m NED elevation,
aggegated to Level III ecoregions. & m/m (slope ratio)\\
TWI Mean & Categorical version of the Topographic Wetness Index (TWI), created by
binning continuous TWI values into ordinal classes (e.g., quartiles or
hydrologically relevant thresholds). Simplifies terrain and moisture
interactions for modeling and interpretability. & unitless index\\
\addlinespace
TWI Modal & Zonal mode of 10 m NED Topographic Wetness Index (TWI) values
aggegated to Level III ecoregions. & unitless index\\
TWI Class & Zonal assignment of ordinal NED Topographic Wetness Index (TWI)
classes, derived from 10 m NED, aggegated to Level III ecoregions. & class label (ordinal ID)\\
Elevation Range & Zonal median elevation range (max minus min elevation) from 10 m NED
elevation data, aggregated to NHDPlusHD catchments. & meters\\
Aspect (Cosine) & Zonal median of the cosine of aspect, derived from 10 m NED elevation
data, aggregated to NHDPlusHD catchments. & unitless [minus 1 to 1]\\
Aspect (Sine) & Zonal median of the sine of aspect, derived from 10 m NED elevation
data, aggregated to NHDPlusHD catchments. & unitless [minus 1 to 1]\\
\addlinespace
Curvature IQR & IQR of combined curvature types, derived from 10 m NED data,
aggregated to NHD+ catchments. & unitless\\
Planform Curvature & Zonal median of planform (horizontal) curvature, derived from 10 m NED
data, aggregated to NHDPlusHD catchments. & unitless\\
Profile Curvature & Zonal median of profile (vertical) curvature, derived from 10 m NED
elevation data, aggregated to NHDPlusHD catchments. & unitless\\
Watershed Elongation & Watershed elongation ratio, calculated from catchment geometry,
aggregated to NHDPlusHD catchments. & unitless ratio (m per m)\\
Watershed Circularity & Watershed circularity ratio, calculated from catchment geometry,
aggregated to NHD+ catchments. & unitless ratio (sq-m per sq-m)\\*
\end{longtable}
\endgroup{}

\vspace{0.5cm}

\begingroup\fontsize{8}{10}\selectfont

\begin{longtable}[t]{lp{8cm}l}
\caption{\label{tab:create_tables}Table B: Topography — Notes and Temporal Scale}\\
\toprule
Short Name & Notes & Scale\\
\midrule
\endfirsthead
\caption[]{Table B: Topography — Notes and Temporal Scale \textit{(continued)}}\\
\toprule
Short Name & Notes & Scale\\
\midrule
\endhead

\endfoot
\bottomrule
\endlastfoot
Station Altitude & Elevation of the gage site above sea level. Influences temperature,
snowpack, and runoff timing & Station\\
Mean Slope & Describes average terrain steepness. Helps distinguish rugged uplands
from flatter plains. & Macroregional\\
Median Slope & Reflects typical terrain slope. More robust to outliers than the mean;
useful for evaluating runoff tendency. & Macroregional\\
Altitude Zone & Represents elevational regime. Influences climate, snow duration, and
runoff seasonality. & Macroregional\\
Mean Slope Distribution & Aggregated mean slope at ecoregion scale. Affects stream power, flow
velocity, and sediment dynamics. & Regional\\
\addlinespace
Slope Skewness & Measures asymmetry in slope distribution. Helps identify dominant
terrain types and slope transitions. & Regional\\
Slope Variability (StDev) & Quantifies slope variability including extremes. Suggests rugged or
heterogeneous terrain. & Regional\\
Slope Variability (IQR) & Measures slope variability using IQR. Robust to outliers; reflects
terrain complexity and infiltration dynamics. & Regional\\
Surface Roughness & Captures fine-scale surface texture. Influences ponding, microrelief,
and infiltration pathways. & Subregional\\
TWI Mean & Indicates average terrain wetness potential. Combines slope and
upstream drainage area. Lower classes indicate drier conditions;
higher classes represent wetter zones with greater runoff accumulation
potential. & Subregional\\
\addlinespace
TWI Modal & Shows most frequent wetness class. Useful for identifying dominant
moisture or runoff conditions. & Subregional\\
TWI Class & Ordinal wetness class derived from TWI. Simplifies terrain–moisture
interactions; lower values = drier, higher = wetter zones. & Subregional\\
Elevation Range & Captures elevation spread within catchment. Linked to slope energy,
erosion risk, and runoff response. & Local\\
Aspect (Cosine) & Quantifies N to S slope orientation. North-facing slopes (values > 0)
retain snow longer and reduce evapotranspiration. & Local\\
Aspect (Sine) & Quantifies E to W slope orientation. East-facing slopes (values > 0)
warm earlier, affecting snowmelt and phenology. & Local\\
\addlinespace
Curvature IQR & Measures variation in terrain curvature. High values imply complex
flow paths and erosion susceptibility. & Local\\
Planform Curvature & Identifies slope convergence/divergence. Negative = concave
(channels), positive = convex (ridges). & Local\\
Profile Curvature & Represents slope change along flow lines. Helps distinguish
depositional vs. erosional terrain zones. & Local\\
Watershed Elongation & Describes catchment shape. More elongated basins have faster runoff
and flashier peaks. & Local\\
Watershed Circularity & Captures catchment compactness. Higher values suggest faster, more
concentrated runoff response. & Local\\*
\end{longtable}
\endgroup{}

\clearpage

\section*{ Watershed Metrics }
\addcontentsline{toc}{section}{ Watershed Metrics }

\begingroup\fontsize{8}{10}\selectfont

\begin{longtable}[t]{lp{8cm}l}
\caption{\label{tab:create_tables}Table A: Watershed Metrics — Definition and Units}\\
\toprule
Short Name & Description & Unit\\
\midrule
\endfirsthead
\caption[]{Table A: Watershed Metrics — Definition and Units \textit{(continued)}}\\
\toprule
Short Name & Description & Unit\\
\midrule
\endhead

\endfoot
\bottomrule
\endlastfoot
Watershed Area & Point-based (gage-level) data used as a station-level predictors, not
derived from zonal aggregation. & watershed area (sq-km)\\
Clay Fraction & Zonal modal fraction of surface soil from zero to 5 cm composed of
clay particles from SSURGO2, aggregated to Level II ecoregions. & proportion [0 to 1]\\
Silt Fraction & Zonal modal fraction of surface soil from zero to 5 cm composed of
silt particles from SSURGO2, aggregated to Level II ecoregions. & proportion [0 to 1]\\
Sand Fraction & Zonal modal fraction of surface soil from zero to 5 cm composed of
sand particles from SSURGO2, aggregated to Level II ecoregions. & proportion [0 to 1]\\
Median Stream Order & Zonal median of Strahler stream order derived from NHDPlus v2.1
flowlines, aggregated to Level II ecoregions. & Strahler order\\
\addlinespace
Max Stream Order & Zonal maximum stream order in the network derived from NHDPlus v2.1
flowlines, aggregated to Level II ecoregions. & Strahler order\\
Soil Permeability & Zonal median of STATSGO2 permeability values aggegated to Level III
ecoregions. & mm/h\\
Soil Runoff Class & Zonal modal soil hydrologic group (A to D) from STATSGO2 aggegated to
Level III ecoregions. & class A to D (ordinal)\\
Flow Accumulation & Zonal median of 10 m NED flow accumulation (upslope contributing area)
aggegated to Level III ecoregions. & contributing area (sq-km)\\
Stream Density & Zonal mean of stream length per catchment area, derived from NHDPlusHD
flowlines, aggregated to NHDPlusHD catchments. & meters per square meter\\
\addlinespace
Mean Flow Length & Zonal median of average flow path length, derived from NHDPlusHD
flowlines, aggregated to NHDPlusHD catchments & meters\\
Stream Slope & Zonal median of stream slope (elevation drop per unit flow length),
derived from NED and NHDPlus data, aggregated to NHDPlusHD catchments. & slope ratio (m/m)\\
Relief Ratio & Relief ratio (elevation range divided by stream length) from 10 m NED
elevation data, aggregated to NHDPlusHD catchments. & m/m\\*
\end{longtable}
\endgroup{}

\vspace{0.5cm}

\begingroup\fontsize{8}{10}\selectfont

\begin{longtable}[t]{lp{8cm}l}
\caption{\label{tab:create_tables}Table B: Watershed Metrics — Notes and Temporal Scale}\\
\toprule
Short Name & Notes & Scale\\
\midrule
\endfirsthead
\caption[]{Table B: Watershed Metrics — Notes and Temporal Scale \textit{(continued)}}\\
\toprule
Short Name & Notes & Scale\\
\midrule
\endhead

\endfoot
\bottomrule
\endlastfoot
Watershed Area & Commonly used as a spatial predictor in hydrologic modeling. & Station\\
Clay Fraction & Indicates surface soil texture. Higher clay content reduces
infiltration and increases runoff potential. & Regional\\
Silt Fraction & Indicates intermediate soil texture. Balances water retention and
infiltration capacity. & Regional\\
Sand Fraction & Represents infiltration rate of surface soils. Low values indicate
higher runoff and steeper hydrographs. & Regional\\
Median Stream Order & Represents typical drainage network complexity. Based on Strahler
order, indicating average stream hierarchy within the watershed. & Regional\\
\addlinespace
Max Stream Order & Indicates the largest stream order within the catchment. Useful for
identifying drainage basin scale and extent. & Regional\\
Soil Permeability & Represents infiltration rate of surface soils. Low values indicate
higher runoff and steeper hydrographs. & Subregional\\
Soil Runoff Class & USDA hydrologic soil group (A to D). Classifies soils by infiltration
capacity and storm response. & Subregional\\
Flow Accumulation & Estimates upslope contributing area. Higher values suggest greater
potential for water accumulation and flooding. & Subregional\\
Stream Density & Ratio of total stream length to catchment area. Higher density
reflects greater drainage dissection and shorter runoff travel times. & Local\\
\addlinespace
Mean Flow Length & Average flow path length to outlet. Longer paths typically indicate
delayed runoff response. & Local\\
Stream Slope & Elevation drop per unit stream length. Influences runoff velocity and
hydrograph shape. & Local\\
Relief Ratio & Ratio of watershed relief to main channel length. High values indicate
steep basins and higher hydrologic energy. & Local\\*
\end{longtable}
\endgroup{}

\clearpage

\section*{ Land Cover }
\addcontentsline{toc}{section}{ Land Cover }

\begingroup\fontsize{8}{10}\selectfont

\begin{longtable}[t]{lp{8cm}l}
\caption{\label{tab:create_tables}Table A: Land Cover — Definition and Units}\\
\toprule
Short Name & Description & Unit\\
\midrule
\endfirsthead
\caption[]{Table A: Land Cover — Definition and Units \textit{(continued)}}\\
\toprule
Short Name & Description & Unit\\
\midrule
\endhead

\endfoot
\bottomrule
\endlastfoot
NLCD Cultivated Fraction & Proportion of NLCD cultivated land: cropland (81), aggregated to
macroregional zones. & proportion [0 to 1]\\
NLCD Forest Fraction & Proportion of NLCD forested land: deciduous (41), evergreen ( 42),
mixed (43), aggregated to macroregional zones. & proportion [0 to 1]\\
NLCD Rangeland Fraction & Proportion of NLCD shrubland (52), grassland (71) and pastureland
(82),, aggregated to macroregional zones. & proportion [0 to 1]\\
NLCD Developed Fraction & Proportion of NLCD developed land open space (21), low-intensity
(22), medium-intensity (23), high-intensity (24), aggregated to
macroregional zones. & proportion [0 to 1]\\
NDVI Amplitude & Zonal median of annual NDVI amplitude for 2016 MODIS data, aggregated
to Level II ecoregions. & unitless [minus 1 to 1]\\
\addlinespace
NDVI IQR & Zonal median of the interquartile range (IQR) of NDVI values for 2016
MODIS data, aggregated to Level II ecoregions. & unitless [minus 1 to 1]\\
Growing Season Length & Zonal median of the length of the 2016 growing season (NDVI-based),
aggregated to Level II ecoregions. & days\\
Peak NDVI & Zonal median of the seasonal peak NDVI for 2016 MODIS data, aggregated
to Level II ecoregions. & unitless [minus 1 to 1]\\
Land Cover \% (MODIS) & Zonal median of percent cover of dominant land class (cropland,
forest, etc.) from MODIS 2016 land cover, aggegated to Level III
ecoregions. & proportion [0 to 1]\\
Land Use Diversity (MODIS) & Normalized Shannon-Weiner Diversity Index of MODIS land classes
aggegated to Level III ecoregions. & unitless index (Shannon-Weiner)\\
\addlinespace
Land Cover \% (NLCD) & Zonal median percent cover by dominant NLCD class aggegated to
NHDPlusHD catchments. & proportion [0 to 1]\\
Land Use Diversity (NLCD) & Normalized Shannon Weiner Diversity Index of NLCD land cover classes,
aggregated to NHDPlusHD catchments. & unitless index (Shannon-Weiner)\\*
\end{longtable}
\endgroup{}

\vspace{0.5cm}

\begingroup\fontsize{8}{10}\selectfont

\begin{longtable}[t]{lp{8cm}l}
\caption{\label{tab:create_tables}Table B: Land Cover — Notes and Temporal Scale}\\
\toprule
Short Name & Notes & Scale\\
\midrule
\endfirsthead
\caption[]{Table B: Land Cover — Notes and Temporal Scale \textit{(continued)}}\\
\toprule
Short Name & Notes & Scale\\
\midrule
\endhead

\endfoot
\bottomrule
\endlastfoot
NLCD Cultivated Fraction & Indicates anthropogenic land use in agricultural regions. Associated
with increased runoff, reduced infiltration, and higher ET demand. & Macroregional\\
NLCD Forest Fraction & Forest cover enhances interception and storage. Typically reduces peak
flows and lowers skew. & Macroregional\\
NLCD Rangeland Fraction & Moderates evapotranspiration and infiltration rates; especially
important in mixed rangeland systems. & Macroregional\\
NLCD Developed Fraction & Urban development increases imperviousness; strongly associated with
peak flow generation and positive flood skew. & Macroregional\\
NDVI Amplitude & Captures seasonal variation in greenness. High amplitude reflects
strong vegetation response to climate fluctuations. & Regional\\
\addlinespace
NDVI IQR & Describes central spread of seasonal greenness. Robust to outliers and
well-suited to fragmented or transitional landscapes. & Regional\\
Growing Season Length & Indicates length of vegetative activity. Influences evapotranspiration
rates and baseflow timing. & Regional\\
Peak NDVI & Represents seasonal greenness maximum. Serves as a proxy for peak
vegetation productivity. & Regional\\
Land Cover \% (MODIS) & Summarizes dominant land cover class by proportion within a zone.
Helps interpret complex or fragmented landscapes by reducing
pixel-level detail to a single representative class. & Subregional\\
Land Use Diversity (MODIS) & Reflects land use heterogeneity. Higher values indicate fragmented or
mixed land cover, which may affect infiltration and runoff patterns. & Subregional\\
\addlinespace
Land Cover \% (NLCD) & Shows dominant land cover proportion at the catchment scale. Useful
for linking terrain features with hydrologic behavior. & Local\\
Land Use Diversity (NLCD) & Measures land cover diversity within a catchment. Supports modeling of
transitions or mosaic land use patterns. & Local\\*
\end{longtable}
\endgroup{}

\clearpage

\end{document}
